\chapter{Metodologia}
%não se esqueça de definir uma label única para utilizar no comando \ref
\label{chap:metodologia}

A metodologia é o coração do seu TCC - ou pelo menos deveria ser. É aqui que você pega tudo o que fez (ou tentou fazer) e descreve como se tivesse seguido um plano super racional, metódico e cuidadosamente pensado. Spoiler: todo mundo sabe que, em muitos casos, foi mais tipo \enquote{tentei várias coisas até funcionar.}

Nessa seção, não se esqueça de descrever o trabalho quanto à natureza, objetivo, abordagem e procedimento. 

\section{Questões de pesquisa}
\label{sub:questoes_pesquisa}

Conforme alinhado na seção \ref{sec:objetivos}:

\begin{itemize}
    \item \textbf{QP1}: Qual o tempo ideal para esperar uma pizza ficar pronta?
    \item \textbf{QP2}: Qual a temperatura ideal para fazer uma pizza?
    \item \textbf{QP3}: Qual a forma ideal de fazer uma pizza?
\end{itemize}

\section{Estratégia de pesquisa}
\label{sub:estrategia_pesquisa}

A busca pelos estudos foi realizada utilizando as bases de dados acadêmicas \textit{\gls{IEEE} Xplore} e \textit{\gls{ACM} Digital Library}, juntamente ao motor de busca \textit{Google Scholar}. A seguir, um exemplo de string de busca: \newline

\textit{(\enquote{survey} OR \enquote{systematic review}) AND (\enquote{adversarial attack} OR \enquote{adversarial defense}) AND (\enquote{computer vision})}\newline

\section{Procedimento de seleção}
\label{sub:selecao_pesquisa}

O procedimento de seleção dos estudos geralmente baseia-se na aplicação de critérios de inclusão e exclusão. Aqui você deve demonstrar o seu rigor científico.

\subsection{Critérios de inclusão}
\label{sub:criterios_inclusao}

Os critérios de inclusão foram definidos conforme elencados na lista:

\begin{itemize}
    \item O estudo aborda a forma ideal de fazer pizzas.
    \item O estudo aborda softwares de cozimento.
\end{itemize}

\subsection{Critérios de exclusão}
\label{sub:criterios_exclusao}

Os critérios de exclusão foram definidos conforme elencados na lista:

\begin{itemize}
    \item O estudo na realidade aborda hambúrgueres.
    \item É um artigo duplicado.
\end{itemize}


%exemplo de código fonte, as configurações estão no arquivo packages.tex
% \lstinputlisting[language=PHP, 
% caption=Texto texto texto texto
% ,label=lst:exemplocodigo1]{chapters/trechos_codigo/funcoescatdinamicas.m}

% \hspace{4cm}
% \hfill
% \begin{minipage}[t]{.65\textwidth}
% \fonteautor
% \end{minipage}


% Texto texto texto texto texto texto texto texto texto texto texto texto texto texto texto texto texto texto texto texto texto texto texto texto texto texto texto texto texto texto texto texto texto texto texto texto, referente ao Código Fonte \ref{lst:exemplocodigo1} função \texttt{nome\_funcao}.

% \subsection{Texto Texto Texto}
% \label{subsec:algumlabel}

% Texto texto texto texto texto texto texto texto texto texto texto texto texto texto texto texto texto texto texto texto texto texto texto texto texto texto texto texto texto texto texto texto texto texto texto texto.

% \section{Texto Texto}
% \label{sec:algumlabel2}
% %exemplo de quadro
% Texto texto texto texto texto texto texto texto texto texto texto texto texto texto texto texto texto texto texto texto texto texto texto texto texto texto texto texto texto texto texto texto texto texto texto texto, ver Quadro \ref{quad:exemplo_de_quadro}.

% \input{chapters/metodologia/quadros/exemplo_de_quadro}


% \subsection{Texto Texto Texto}
% \label{subsec:algumlabel2}
% Texto texto texto texto texto texto texto texto texto texto texto texto texto texto texto texto texto texto texto texto texto texto texto texto texto texto texto texto texto texto texto texto texto texto texto texto.

% \subsection{Texto Texto Texto}
% \label{subsec:algumlabel3}

% Texto texto texto texto texto texto texto texto texto texto texto texto texto texto texto texto texto texto texto texto texto texto texto texto texto texto texto texto texto texto texto texto texto texto texto texto.
